% !TEX root = ../main.tex

\clearpage
% Source PDF page 5.
\section{第二章\quad 奶龙的诞生与发展史}

\subsection{一、原版奶龙的诞生}

2017 年，第七印象文化传媒（深圳）有限公司成立，命运的齿轮开始转动。

2019 年底，原创奶龙 IP 开始创作。

2020 年 1 月 22 日，奶龙 IP 在短视频平台正式上线；同年抖音粉丝破千万，单条视频 3 天点播破亿。2021 年，《奶龙》短视频合集在优酷、爱奇艺、腾讯同步更新；有声读物专辑上线，8 个月播放量破亿。

2023 年，第一季番剧《奶龙搞怪大作战》播出，至此奶龙进入了高中生和大学生的视野。

2024 年，第二季番剧《奶龙与小七之大战暴暴龙》播出，奶龙开始在高中生和大学生之间爆火。

2025 年 1 月，第三季番剧《奶龙之奇妙冒险》上线，一个月内全网播放量破 1 亿，热度达到了前所未有的顶峰。

奶龙作为新一代子供向作品，无疑是极其成功的，不仅在国内大受欢迎，影响力甚至辐射到了东南亚地区。

然而对于大多数青年来说，奶龙的形象可以说是一言难尽。

\subsection{二、原版奶龙在传播中的重要地位}

原版奶龙在大多数大学生心目中，毫无疑问是一个糖完了的形象。~(见\figref{fig:02-1})

% !TEX root = ../../main.tex
\begin{figure}[H]
  \centering
  \includegraphics[width=73.27bp,height=0.76\textheight,keepaspectratio]{figures/chapter02/figure02_1.pdf}
  \renewcommand{\thefigure}{2.1}
  \caption{被 P 图恶搞的奶龙}
  \label{fig:02-1}
\end{figure}


原版奶龙最受欢迎的台段毫无疑问是“我是奶龙，我才是奶龙”。不过并不是因为这一段好看，而是因为这一段很糖。

“我觉得真假奶龙这一段简直是智斗天花板，当小七看着蹦出来的两个奶龙时，冷汗瞬间就从我的脑袋上冒下来了，在假奶龙说出我会喷火后，我全身的血液都在上涌，我的思维推理已经跟不上剧情的演绎了，当最后真奶龙说出那句你才是奶龙的时候，我早已浑身湿透了，我当场就被这完美的智斗剧情折服了。”

这种真假美猴王的桥段大家早已看过很多次了，而奶龙自带的糖气和鲱鱼罐头引发的“你才是奶龙”又让这一段糖上加糖。

在倒放大法的作用下，“我是奶龙，我才是奶龙”又变成“我才是奶龙，欧蓝的手”（找不到原视频了 qwq），这一妙手偶得之（并非偶得）让奶龙的糖度爆表，一度为人称道。

随后，奶龙与贝利亚进行了深度绑定，在抽象道路上一去不复返。

虽然原版奶龙本身很正经，也受到少数青年的喜爱，但是却给大多数大学生意外提供了糖度，毫无疑问为接下来奶蛙的登场打好了基础。

